%=============================================================================
\section{完整工作流总结}
%=============================================================================

\begin{enumerate}
  \item \textbf{本地准备}：开发代码，准备好 \texttt{train.py}、配置文件和 SLURM 脚本

  \item \textbf{上传到集群}：
\begin{lstlisting}[language=bash]
rsync -azP --exclude='.venv' --exclude='__pycache__' \
  ./my_project/ your_utln@login-prod.pax.tufts.edu:~/project/
\end{lstlisting}

  \item \textbf{SSH 登录，设置环境}：
\begin{lstlisting}[language=bash]
ssh your_utln@login-prod.pax.tufts.edu
module load miniforge/25.3.0 && module load cuda/12.9.0
conda activate rl_env
\end{lstlisting}

  \item \textbf{交互式调试}——先跑一个任务确认一切正常：
\begin{lstlisting}[language=bash]
srun -p gpu -t 0-1:00:00 --gres=gpu:1 --mem=8g --cpus-per-task=4 --pty bash
python train.py --env CartPole-v1 --arch PPO --seed 0 --output-dir test_run
exit
\end{lstlisting}

  \item \textbf{用 seff 检查资源使用}，调整请求量

  \item \textbf{批量提交所有实验}：
\begin{lstlisting}[language=bash]
bash scripts/submit_all.sh
\end{lstlisting}

  \item \textbf{监控进度}：
\begin{lstlisting}[language=bash]
squeue --me                           # 查看状态
watch -n 60 'squeue --me | wc -l'     # 持续监控剩余任务数
\end{lstlisting}

  \item \textbf{下载结果}：
\begin{lstlisting}[language=bash]
# 在本地运行
rsync -azP your_utln@login-prod.pax.tufts.edu:~/project/results/ ./results/
python collect_results.py
\end{lstlisting}
\end{enumerate}
